% SOURCE_AUTHORITY: sga5_fr_workpass.tex, lines 8258--8286 (LF-normalized unit).
% SOURCE_SHA256: 791F4EFFC5E02832D5D77ED03518C8156D6F07E4C8238B03545DB93D883FBB28
\subsubsection*{1.1.4}
Sea \(F=(F_n)_{n\in\mathbb N}\) un haz \(\ell\)-ádico constructible. Como el anillo \(\mathbb Z_\ell\) actúa de manera evidente sobre cada haz \(F_n\), actúa asimismo sobre \(F\). Además, si \(F\) y \(G\) son dos haces \(\ell\)-ádicos constructibles y
\[
u:F\longrightarrow G
\]
es un morfismo, \(u\) es compatible con las acciones de \(\mathbb Z_\ell\) sobre \(F\) y \(G\), respectivamente. Esto muestra que \(\ell\text{-}\adc(X)\) está dotada canónicamente de una estructura de \(\mathbb Z_\ell\)-categoría abeliana.

Esto justifica la terminología siguiente:

\begin{statement}{Terminología}
En lo sucesivo, los haces \(\ell\)-ádicos constructibles se llamarán \(\mathbb Z_\ell\)-haces constructibles.
\end{statement}

Conviene no confundir esta noción con la de \((\mathbb Z_\ell)_X\)-módulo, que no se utilizará en lo que sigue.

La categoría \(\ell\text{-}\adc(X)\) se denotará asimismo por
\[
\mathbb Z_\ell\text{-}\fc(X).
\]

\subsubsection*{1.2. \(\mathbb Z_\ell\)-haces constantes torcidos constructibles}

\begin{statement}{Definición 1.2.1}
Un \(\mathbb Z_\ell\)-haz constructible
\[
F=(F_n)_{n\in\mathbb N}
\]
se dice constante torcido si y solo si los haces \(F_n\) son localmente constantes sobre \(X_{\mathrm{\acute et}}\).
\end{statement}
